% ======================================================================
\section{Introduction}\label{sec:intro}
% ======================================================================

At high energy the wave function of a hadron or nucleus contains a large number of soft gluons, and in a sufficiently dense system their density saturates~\cite{Gribov:1983ivg,Mueller:1985wy}. A description of this regime from first principles was proposed by McLerran and Venugopalan~\cite{McLerran:1993ni,McLerran:1993ka}: the fast (valence) partons of a large nucleus are treated as static sources of colour charge, with a density $\rho^a(\v x)$ in the transverse plane; the soft gluons are the classical Yang--Mills field radiated by these sources; and observables are obtained by averaging over the sources with a Gaussian weight. This picture has developed into the effective theory known as the Color Glass Condensate (CGC)~\cite{Iancu:2003xm,Weigert:2005us,Gelis:2010nm,Kovchegov:2012mbw}. The dependence of observables on the rapidity interval over which the soft modes are integrated out is governed by the BFKL equation~\cite{Kuraev:1977fs,Balitsky:1978ic} in the dilute regime and by the non-linear Balitsky--Kovchegov (BK) and JIMWLK equations~\cite{Balitsky:1995ub,Kovchegov:1999yj,Kovchegov:1999ua,Jalilian-Marian:1997qno,Jalilian-Marian:1997jx,Jalilian-Marian:1997ubg,Kovner:2000pt,Weigert:2000gi,Iancu:2000hn,Iancu:2001ad,Ferreiro:2001qy} in the dense regime, which replace the Gaussian weight by an evolving distribution of the sources.

The simplest observable that is sensitive both to the projectile and to the target wave functions is the single inclusive gluon spectrum in the collision of a dilute projectile with a dense target. At leading order (LO), for a gluon produced at mid rapidity, i.e.\ far in rapidity from both the projectile and the target, the spectrum is given by the formula of Kovchegov and Tuchin~\cite{Kovchegov:2001sc,Kovchegov:1998bi}: the gluon is emitted from the valence charges of the projectile in the Weizs\"acker--Williams field, and the emission amplitude, together with the colour rotation of the emitting charge, scatters eikonally on the target. The cross section is a product of the projectile colour charge correlator and of a target operator built from two adjoint Wilson lines; in the limit of a dilute target it reduces to the $k_T$-factorised form with the Lipatov vertex~\cite{Gribov:1983ivg,Kharzeev:2003wz,Blaizot:2004wu}. This is a result of the classical theory: it is quadratic in the charges of the projectile, and after the average over a colour-neutral projectile it does not depend on whether these charges are numbers or operators. Its leading logarithmic (LL) rapidity evolution was obtained by Kovner and Lublinsky~\cite{Kovner:2006wr} (see also~\cite{Kovner:2006ge}) and, within the Schwinger--Keldysh formulation, by Gelis, Lappi and Venugopalan~\cite{Gelis:2008rw}: the projectile is to be evolved from its own rapidity down to the rapidity of the produced gluon, and the target from the rapidity of the produced gluon to its own rapidity.

Beyond leading order the situation is different, and the picture of classical fields of static, commuting colour charges with a Gaussian weight, which was designed for the leading-order problem and solves it, is not by itself a complete description of gluon production. Three elements have to be added to it. First, the colour charges of the valence partons are operators, $[\rho^a(\v x),\rho^b(\v y)]=if^{abc}\rho^c(\v x)\delta^{(2)}(\v x-\v y)$~\cite{Kovner:2005pe,Kovner:2005nq,Kovner:2005jc,Kovner:2005en}. At the order $g^4$ of the next-to-leading order (NLO) correction the amplitudes contain products of three and four charges, and a commutator of two charges is of the same order as the two-charge terms that carry the evolution and the NLO coefficient. The order of the charges therefore matters: the commutators leave contact terms, and the totally symmetrised product of three charges brings in the three-point function of the projectile with the symmetric structure constants $d^{abc}$, which is odd under charge conjugation and absent for a Gaussian weight. Second, the logarithms of an NLO calculation are of several kinds: those of the rapidity evolution of the projectile and of the target, the collinear logarithms of the fragmentation of the produced gluon, which are governed by the DGLAP equations~\cite{Gribov:1972ri,Altarelli:1977zs,Dokshitzer:1977sg}, and the ultraviolet logarithms that make the coupling run. In the leading-order framework they are supplied by separate considerations; in a complete description all of them, with their scales and their overlaps, have to emerge from one calculation. Third, the terms of relative order $\alpha_s$ that are not enhanced by any logarithm, which are needed for any quantitative use of the leading-order formula, have not been available for this observable.

Many of the ingredients have become available over the last years. The evolution equations themselves are known at NLO: the NLO BK equation~\cite{Balitsky:2007feb} and the NLO JIMWLK Hamiltonian~\cite{Balitsky:2013fea,Kovner:2013ona,Kovner:2014lca,Grabovsky:2013mba}, together with the running-coupling corrections to the LO kernels~\cite{Balitsky:2006wa,Kovchegov:2006vj,Kovchegov:2006wf}. In particular, Lublinsky and Mulian~\cite{Lublinsky:2016meo} computed the soft gluon light-cone wave function of a dilute projectile to order $g^3$ and derived the NLO JIMWLK Hamiltonian from it. For particle production, the NLO corrections to single inclusive hadron production at forward rapidity were computed in the hybrid formalism~\cite{Dumitru:2005gt,Altinoluk:2011qy} by Chirilli, Xiao and Yuan~\cite{Chirilli:2011km,Chirilli:2012jd}, and the structure of the rapidity subtraction and of the kinematic constraints in this framework was analysed in~\cite{Stasto:2013cha,Altinoluk:2014eka,Ducloue:2016shw}. Running-coupling corrections to inclusive gluon production at mid rapidity were obtained by Horowitz and Kovchegov~\cite{Horowitz:2010yg}, and the relation between $\beta_0$-dependent terms of NLO calculations and the running of the coupling was clarified in~\cite{Kovner:2023vsy}. The general structure of the NLO correction to inclusive spectra in the CGC was discussed by Gelis and Laidet~\cite{Gelis:2012ri}. Corrections of higher order in the projectile colour charge (``first saturation corrections'' in the projectile) were computed in the classical approximation in~\cite{Chirilli:2015tea,Li:2021zmf,Li:2021ntm}, and their contribution to single inclusive production, odd in $\v k$, was analysed in~\cite{McLerran:2016snu,Kovchegov:2018jun}. A complete NLO calculation of the single inclusive gluon spectrum at mid rapidity has, to our knowledge, not been presented.

The purpose of this paper is to supply the missing pieces for this observable. In a previous paper~\cite{Radhakrishnan:2026soft} the soft gluon wave function of a dilute projectile and the associated unitary evolution (dressing) operator $\Omega$ were constructed up to order $g^2$, in pure Yang--Mills theory, within light-cone perturbation theory (LCPT)~\cite{Lepage:1980fj,Brodsky:1997de}. The operator $\Omega$ was written as a normal-ordered series in soft gluon creation and annihilation operators with coefficients that are functionals of the non-commuting valence charge density $\rho$, and the coefficients were fixed by matching to LCPT and by unitarity. The NLO cross section requires in addition the coefficient of $\Omega$ that creates one gluon at order $g^3$; its ingredients, the $O(g^3)$ one-gluon components of the light-cone wave function, were computed by Lublinsky and Mulian~\cite{Lublinsky:2016meo}, and we recompute them here with the longitudinal cut-offs and the operator ordering required by the cross section. The observable is the expectation value of the number operator of dressed gluons in the scattered, dressed projectile state~\cite{Kovner:2006wr}, so that the full calculation reduces to an algebraic evaluation of products of the coefficients of $\Omega$ and of Wilson lines, in which the order of the charges is kept, followed by longitudinal and transverse integrations. No average over the projectile is taken until the end, and no model for it is assumed.

Our main results are the following.
\begin{itemize}
\item The $O(g^4)$ cross section naturally decomposes into three groups of terms: interference of the LO amplitude with the $O(g^3)$ one-gluon amplitude (Group~I); the squares of the $O(g^2)$ two-gluon amplitudes, with one of the gluons integrated over (Group~II); and the interference terms between different two-gluon amplitudes (Group~III). We list all products of coefficients that contribute, with their longitudinal integration ranges and their order in $\rho$ (Section~\ref{sec:observable}).
\item All terms enhanced by $\log(\vee/\Lambda)$, where $\Lambda<p^+<\vee$ is the longitudinal momentum range of the soft modes, are the LL evolution of the LO result. The logarithm splits at the longitudinal momentum $k^+$ of the produced gluon into $\log(\vee/k^+)+\log(k^+/\Lambda)$; the coefficient of $\log(\vee/k^+)$ is exactly the BFKL evolution of the projectile charge correlator inserted into the LO formula, and the coefficient of $\log(k^+/\Lambda)$ is exactly the JIMWLK Hamiltonian acting on the target operator of the LO formula. We give both terms for an arbitrary (ordered) projectile correlator and in colour-singlet form. The contributions quartic in $\rho$, all of which carry $\log(\vee/\Lambda)$, cancel identically for commuting charges, and the totally symmetrised cubic terms cancel as well; with non-commuting charges, their reordering produces the two-$\rho$ terms that complete the evolution (Section~\ref{sec:rapidity}).
\item The remaining terms are a collinear term, associated with the fragmentation of the produced gluon and governed by the gluon--gluon splitting function, a running-coupling term, and finite terms. We give an account of all poles $1/\epsilon$: the ultraviolet poles of the virtual corrections to the produced gluon cancel among themselves, the coupling is renormalised by the normalisation of the dressed produced gluon, and the collinear poles combine into the full $\MS$ splitting function $\mathcal P_{gg}$ and are absorbed into the fragmentation function. The running-coupling term is the $\frac{\beta_0}{2}\delta(1-\xi)$ part of the collinear term; it sets the scale of the LO coupling. The scheme is conventional dimensional regularisation with the $\MS$ coupling and the standard $\MS$ number-density fragmentation function, with which standard one-hadron-inclusive coefficient functions are reproduced (Section~\ref{sec:dglap_rc}).
\item We give the finite terms Group by Group and Row by Row, after the cancellations that are exact, in two forms: for an arbitrary projectile, with every product of three charges written as its totally symmetrised product plus two-charge contact terms, and averaged over the projectile, in colour-singlet form. All colour structures are written, as at leading order, through differences of Wilson lines, and all longitudinal integrals with plus prescriptions in a momentum fraction. After the average the terms cubic in $\rho$ do not vanish: their symmetrised part, odd in $\v k$, is proportional to the symmetrised ($d^{abc}$) three-point function of the projectile, which is kept general, and to C-odd structures of the target, while the commutators of the ordered charges leave contact terms, even in $\v k$, which are present for any projectile. The large-distance tails of individual terms cancel in pairs of a virtual and a real term, without any condition on the target at large distance; we state which terms have to be integrated together. A compact summary of the result is given in Section~\ref{sec:result_summary} (Section~\ref{sec:finite}, Appendices~\ref{app:finite} and~\ref{app:tails}).
\item In the dilute--dilute limit the real emission reproduces exactly the two-gluon production cross section in quasi-multi-Regge kinematics~\cite{Fadin:1998py,Bartels:2006hg}; for a point-like projectile the coefficients of the rapidity logarithms in the virtual correction are the gluon Regge trajectories. The comparison of the finite part of the virtual correction with the one-loop central-emission vertex, which requires the real emission with matched regulators, remains open, as do the other points listed in Section~\ref{sec:openpoints} (Section~\ref{sec:checks}).
\end{itemize}
The result can be written schematically as
\begin{align}\label{eq:schematic}
\frac{d^3N}{d^3k} &= \frac{d^3N_{\rm LO}}{d^3k}
+ \log\frac{\vee}{k^+}\,\frac{d^3N_{\rm BFKL}}{d^3k}
+ \log\frac{k^+}{\Lambda}\,\frac{d^3N_{\rm JIMWLK}}{d^3k}\notag\\
&\quad+ \frac{d^3N_{\rm DGLAP}}{d^3k}
+ \frac{d^3N_{\rm RC}}{d^3k}
+ \frac{d^3N_{\rm finite}}{d^3k}\,,
\end{align}
where every term on the right-hand side except the first is of relative order $\alpha_s$. In this formula the evolution of the projectile, the evolution of the target, the collinear evolution of the final state, the running of the coupling and the finite terms, with the charges of the projectile treated as operators throughout, come out of a single calculation. In this sense the present work completes the NLO description of the single inclusive gluon spectrum at mid rapidity within the eikonal, dilute-projectile approximation and in pure Yang--Mills theory, up to the open points of Section~\ref{sec:openpoints}; it does not address the dense--dense problem, corrections beyond the eikonal approximation, or quarks.

\subsection*{Roadmap and reading guide}\label{sec:roadmap}

The paper is written so that it can be read at three levels.
\begin{itemize}
\item \emph{The result only.} Section~\ref{sec:notation} (conventions and vocabulary), then Section~\ref{sec:result_summary} with Table~\ref{tab:result} and the reading guide of Section~\ref{sec:howtoread}, and the discussion of Section~\ref{sec:discussion}. The open points are listed in Section~\ref{sec:openpoints}.
\item \emph{The logic of the calculation.} In addition Sections~\ref{sec:setup} and~\ref{sec:observable} (construction of $\Omega$, the observable, leading order, Groups and rows) and the three worked examples of Section~\ref{sec:examples}, in each of which one row is followed from the coefficients of $\Omega$ to its final form. Each section begins with a paragraph that says what is done and why, and ends with a boxed summary.
\item \emph{The details.} Sections~\ref{sec:rapidity}--\ref{sec:finite} and the appendices, where the points that require care are each treated in a paragraph or subsection of their own. Table~\ref{tab:nuances} below lists them with their locations.
\end{itemize}
The sections are organised as follows. Section~\ref{sec:notation} gives the map of the calculation, the table of symbols, the conventions and the vocabulary. In Section~\ref{sec:setup} we summarise the setup: light-cone conventions, the Hamiltonian of the soft modes in the eikonal approximation, light-cone perturbation theory with operator-valued matrix elements, the construction of $\Omega$ order by order, and its coefficients in coordinate space. In Section~\ref{sec:observable} we define the observable, derive the LO cross section step by step, and organise the NLO contributions into Groups and rows. Section~\ref{sec:examples} contains the worked examples. Section~\ref{sec:rapidity} is devoted to the rapidity logarithms and their identification with the BFKL and JIMWLK evolution. The poles, the collinear and the running-coupling terms are discussed in Section~\ref{sec:dglap_rc}, the finite terms in Section~\ref{sec:finite}, and limits, checks and comparisons with the literature in Section~\ref{sec:checks}, which begins with a table of all checks and of what each of them does and does not test. Section~\ref{sec:discussion} discusses what is new, how the result is to be used and what a numerical evaluation requires. We conclude in Section~\ref{sec:conclusions}, where the open points are listed. The appendices collect the Hamiltonian and its matrix elements (Appendix~\ref{app:hamiltonian}), the light-cone wave functions from which the coefficients of $\Omega$ are obtained (Appendix~\ref{app:wavefunctions}), the unitarity relations among the coefficients (Appendix~\ref{app:unitarity}), the Fourier transforms (Appendix~\ref{app:fourier}), the prescriptions and distributions used in the longitudinal and transverse integrals (Appendix~\ref{app:distributions}), the finite terms for an arbitrary projectile (Appendix~\ref{app:finite}) and the analysis of their large-distance behaviour (Appendix~\ref{app:tails}).

\begin{table}[t]
\centering
\renewcommand{\arraystretch}{1.2}
\footnotesize
\begin{tabular}{@{}>{\raggedright\arraybackslash}p{0.60\textwidth}>{\raggedright\arraybackslash}p{0.34\textwidth}@{}}
\toprule
point that requires care & where it is treated\\
\midrule
non-commuting charges: order of operators, the term $\mf A^{(1)}\mf N_2$, the commutator $[\mf N_2,\mf A^{(1)}]$ & Secs.~\ref{sec:lcpt}, \ref{sec:whyA3}, \ref{sec:NLOamplitudes}; App.~\ref{app:unitarity}\\
symmetrisation, contact terms, one-$\rho$ terms & Secs.~\ref{sec:fourrho}, \ref{sec:fin_sym}\\
four-$\rho$ terms and the coherent state & Sec.~\ref{sec:fourrho}\\
which longitudinal region gives which logarithm; kinematics fixes the limits & Secs.~\ref{sec:loglong}, \ref{sec:groups}\\
plus prescriptions, recoil phases, soft remainders, double logarithms & Secs.~\ref{sec:fin_conv}, \ref{sec:nuance_soft}; App.~\ref{sec:fin_soft}\\
pairing of Group~I and Group~III terms; the ``problematic'' term of Row~I.III and the order of the $\v w$ and $\v z$ integrals & Secs.~\ref{sec:fin_together}, \ref{sec:nuance_soft}; App.~\ref{sec:fin_GIIII}\\
poles: ultraviolet versus collinear, the $\frac{11}{3}$ cancellation, the scaleless rule, order of limits & Sec.~\ref{sec:dg_poles}\\
DGLAP: collinear region, splitting function, two end-point subtractions, scheme & Secs.~\ref{sec:dg_origin}, \ref{sec:dg_msbar}, \ref{sec:dg_final}\\
running coupling: what the $\beta_0$ term is, $\mu_R$ and $\mu_F$, which $\frac{11}{6}$ is which, result versus grouping & Secs.~\ref{sec:dg_poles}, \ref{sec:dg_final}, \ref{sec:lit_constant}\\
large-distance tail of Row~II.IV, cancellation $1:1:-2$ & Sec.~\ref{sec:dg_tail}\\
principal value in $\mf B_3$; instantaneous terms as distributions & Secs.~\ref{sec:fin_conv}, \ref{sec:nuance_pv}; App.~\ref{sec:fin_inst}\\
integrability and power counting; Row~I.II kernel; Row~I.V collinear point & Sec.~\ref{sec:fin_together}; App.~\ref{sec:fin_integrable}\\
projectile average: assumptions, $n$ and $\sigma$, parity, C-parity & Secs.~\ref{sec:fin_singlet}, \ref{sec:threerho}\\
the checks and what each does and does not test & Sec.~\ref{sec:checks_table}\\
\bottomrule
\end{tabular}
\caption{The points of the calculation that require care, and where each is discussed.}
\label{tab:nuances}
\end{table}
